% Source brief: government generative-AI chatbots in Australian public schooling.
% Two case initiatives, South Australia (EdChat) and Queensland (Cerego trial,
% then Corella). Prose is citation-ready against References/References.bib.

\section{South Australia: EdChat}
\label{sec:ai:edchat}

\paragraph{Location.}
South Australian government schools, administered statewide by the Department for
Education \autocite{educationSAEdChat2025}.

\paragraph{Overview.}
EdChat is a department-built generative-AI chatbot, developed with Microsoft and
delivery partner Insight, and powered by the Azure OpenAI Service running the same
model family as ChatGPT \autocite{microsoftEdChatContentSafety2024}. Its aim is to
give students and staff a sandboxed conversational assistant that supports learning
and reduces teacher administrative load while suppressing the data-leakage and
harmful-content risks of consumer chatbots \autocite{educationSAEdChat2025}. An
initial eight-week pilot, 1,500 students and 150 educators across eight secondary
schools, completed in August 2023; subsequent proof-of-concept phases involved more
than 10,000 students \autocite{microsoftEdChatContentSafety2024}. Access opened to
all department staff, including principals, teachers, and preschool staff, in late
2024, with a phased rollout to public secondary students from Term 4, 2025
\autocite{innovationausEdChatStatewide2025}. The organising body is the South
Australian Department for Education, under Chief Executive Martin Westwell.

\paragraph{Curriculum.}
EdChat is positioned as a cross-subject learning aid rather than a discrete course.
Documented student uses include practising foreign languages, working mathematics
problems, rephrasing difficult instructions, and probing complex concepts to aid
comprehension and creativity \autocite{microsoftLearningAIEra2024}. Teachers govern
classroom use through a traffic-light convention, red where AI would compromise
assessment of a student's own writing, green where it raises efficiency in research
and planning, integrating the tool into pedagogy rather than displacing it
\autocite{educationSAEdChat2025}.

\paragraph{Infrastructure and scalability.}
The service is hosted in the department's private Azure tenancy with all data stored
in Australia and never transmitted to public ChatGPT servers
\autocite{educationSAEdChat2025}. Azure AI Content Safety screens prompts and
responses for harmful material, and comprehensive interaction logging gives school
staff oversight of student usage, including outside school hours, with filtered-content
alerts routed to principals the next school day rather than monitored in real time
\autocite{microsoftEdChatContentSafety2024}. This managed, single-tenant architecture
underpinned the move from an eight-school pilot to statewide secondary deployment.

\paragraph{Impact.}
During the trial, active student use rose from roughly 20 per cent at the outset to a
peak by the end of the eight weeks as familiarity grew
\autocite{microsoftEdChatContentSafety2024}. Across the wider rollout, 41 per cent of
students and 36 per cent of educators used the tool at least once, with 93 per cent of
in-hours student prompts curriculum-related. Educators applied it to lesson planning,
correspondence, and assessment-support tasks, the stated mechanism for reducing
workload \autocite{microsoftLearningAIEra2024}.

\paragraph{Evaluation and analysis.}
An external evaluation by Telstra-Quantium and The Learning Future (2024) tracked the
10,000-plus trial cohort and found 94 per cent of student interactions related to
school curriculum subjects, the principal evidence cited for the safety and
appropriateness judgement that authorised statewide release
\autocite{educationSAEdChat2025}. The evaluation emphasises curriculum-focused use,
strong cross-subject uptake, and administrative time savings. Independent commentary
notes the absence of published longitudinal learning-outcome data and cautions against
over-reliance, marking the chief gap between demonstrated engagement and demonstrated
attainment effects \autocite{innovationausEdChatStatewide2025}.

\section{Queensland: Cerego Trial and Corella}
\label{sec:ai:queensland}

\paragraph{Location.}
Queensland state schools, administered by the Department of Education, Queensland
\autocite{educationQueenslandGenerativeAI2026}.

\paragraph{Overview.}
Queensland has run two distinct initiatives. First, a 2023 trial, announced 10 October
2023 by Education Minister Grace Grace, deployed a platform from US vendor Cerego to
500 students across ten state schools \autocite{acsQueenslandTrialAI2023}. Second,
Corella, a department-built generative-AI assistant developed with the Department of
Customer Services, Open Data and Small and Family Business, moved from a 15-school
trial toward statewide rollout to every state high school by June 2026
\autocite{educationdailyCorellaRollout2025}. The shared aim is to support teaching and
learning, prepare students for future work, and reduce teacher workload, with the
department framing AI as a supportive tool and not a replacement for educators
\autocite{queenslandAITrialSchools2023}.

\paragraph{Curriculum.}
The Cerego trial covered English, physics, science, health, humanities, and
accounting, delivering adaptive quiz-based learning that adjusts to each student and
draws only on a digitised version of the approved Queensland state-school curriculum,
a design choice that bounds hallucination and protects academic integrity
\autocite{acsQueenslandTrialAI2023}. Corella is a broader text-based assistant for
generating ideas, explanations, summaries, and drafts, used by students for research,
drafting, brainstorming, and revision, and explicitly to teach responsible AI use
\autocite{educationdailyCorellaRollout2025}.

\paragraph{Infrastructure and scalability.}
Cerego is a third-party platform restricted to approved curriculum content
\autocite{acsQueenslandTrialAI2023}. Corella is a large-language-model assistant hosted
within a secure Queensland Government environment and purpose-built for classroom use
rather than relying on public AI services, the basis for scaling from 15 trial schools
to the full state secondary system \autocite{educationQueenslandCorellaTerms2026}. Both
initiatives are governed by the six principles of the Australian Framework for
Generative Artificial Intelligence in Schools, covering privacy, security and safety,
wellbeing, teaching and learning, transparency, and fairness and accountability
\autocite{queenslandAITrialSchools2023}.

\paragraph{Impact.}
The Corella trial across 15 schools, including Mt Gravatt State High in Brisbane,
indicated potential to cut teacher workload by about 25 per cent by automating lesson
planning, assessment, reporting, and data entry, alongside reported gains in student
digital literacy and teacher work-life balance
\autocite{educationdailyCorellaRollout2025}. Statewide access will extend to teachers,
teacher aides, school leaders, and students in Years 7 to 10, with parental consent
required for student use \autocite{educationdailyCorellaRollout2025}.

\paragraph{Evaluation and analysis.}
Evidence to date is trial-stage and largely workload-oriented rather than
attainment-oriented. The headline 25 per cent workload figure derives from the
15-school Corella trial and is reported by government and secondary outlets rather than
a published independent evaluation \autocite{educationdailyCorellaRollout2025}. The
2023 Cerego trial was structured as a controlled exploration of adaptive learning
against the digitised curriculum, with the department stressing accuracy, privacy, and
teacher oversight as evaluation criteria \autocite{queenslandAITrialSchools2023}. As
with EdChat, the principal analytical gap is the lack of public measurement linking
chatbot use to student learning outcomes.
